\uuid{m80O}
\titre{Matrice d'une application linéaire}
\auteur{hanani ; mzari}
\organisation{exo7}
\licence{CC BY-NC-SA 4.0}
\datecreate{2020-01-20}
\niveau{L1}
\module{Algèbre}
\chapitre{Matrice}
\sousChapitre{Matrice et application linéaire}
\difficulte{1}
\theme{Matrice d'une application linéaire}
\exo7qcmid{lille-2.matrice-application.001}

\contenu{
\begin{question}
\qclassification{108.03}
\qtags{motcle=Matrice d'une application linéaire}
On considère $\Rr$ et $\Rr^2$ munis de leurs bases canoniques et $f$ l'application linéaire définie par :
$$\begin{array}{rccc}f:&\Rr^2 &\to& \Rr\\
& (x,y)&\to &y-x.  \end{array}$$
La matrice de $f$ relativement aux bases canoniques est :
\begin{answers}
\good{$\left(\begin{array}{rc}
-1&1\\
\end{array}\right).$}
\bad{$\left(\begin{array}{r}
-1\\
1\\
\end{array}\right).$}
\bad{$\left(\begin{array}{rc}
-1&1\\
0&0\\
\end{array}\right).$}
\bad{$\left(\begin{array}{rcc}
-1&0\\
1&0\end{array}\right).$}
\end{answers}
\begin{explanations} Soit ${\cal {B}}=\{e_1,e_2\}$ et ${\cal {B}}'=\{1\}$ les bases canoniques de $\Rr^2$ et $\Rr$
respectivement. La matrice de $f$ relativement à ces bases est la matrice dont la $1$ère colonne est $f(e_1)=-1$ et la 2ème colonne est $f(e_2)=1$. Cette matrice est : $\left(\begin{array}{rc}
-1&1\\ \end{array}\right).$
\end{explanations}
\end{question}
}
